% GENERATED FILE. Edit canonical lessons, then run npm run generate:books.
% canonical-source-hash: fnv1a64:2943d56adb7b3c06
% canonical-lessons: ZH-C421-gonglu, ZH-C421-gaosu, ZH-W421-su-speed, ZH-C421-chesu, ZH-C421-shoufeizhan, ZH-C421-chaosu

\chapter{On the Expressway}
\label{ch:zh-on-the-expressway}
\hlchaptermodality{421}

\begin{chapteropening}
\textbf{By the end of this chapter:} \emph{I can talk about main roads and the expressway, the speed of a car, toll gates and speeding.}

Complete the last lesson of chapter 421: I can talk about main roads and the expressway, the speed of a car, toll gates and speeding.
\end{chapteropening}

\section[\texorpdfstring{\zh{公路} (gōnglù)}{公路 (gōnglù)}]{\zh{公路} (gōnglù) — highway; main road}

\label{lesson:ZH-C421-gonglu}



\begin{quote}
Before the new one: say the Mandarin for to reverse the car, then the Mandarin for to break the traffic rules.
\end{quote}

\subsection*{You'll want to know: \zh{公路}}
\textbf{\zh{公路}} — \emph{gōnglù} — \textquotedblleft{}highway; main road\textquotedblright{}.

\begin{quote}
\emph{zhè tiáo gōnglù hěn cháng} — This main road is very long.
\end{quote}

\begin{cousinweb}
\textbf{\zh{公路}} (\emph{gōnglù}) is built from \zh{公} 'public' and \zh{路} 'a road', both inherited from Middle Chinese, \emph{kuwng} and \emph{luH} (Guangyun, in Baxter's transcription).
\end{cousinweb}

\begin{grammarlens}[title={What you've built}]
One more word for getting around.
\end{grammarlens}

\subsection*{Your turn}
\begin{itemize}
\raggedright
  \item \emph{Say it:} \emph{gōnglù}
  \item \emph{Say it:} \emph{gōnglù}, once more
  \item \emph{Say it:} \emph{zhè tiáo gōnglù hěn cháng} — This main road is very long.
  \item \emph{Recall:} say the Mandarin for to reverse the car, then the Mandarin for to break the traffic rules, then say \emph{gōnglù} again
\end{itemize}

\begin{tcolorbox}[breakable,colback=teal!4,colframe=teal!35!black,title={Before you move on}]
What does \zh{公路} mean? (\textquotedblleft{}Highway; main road\textquotedblright{}.) Say it once more.
\end{tcolorbox}
\section[\texorpdfstring{\zh{高速} (gāosù)}{高速 (gāosù)}]{\zh{高速} (gāosù) — expressway}

\label{lesson:ZH-C421-gaosu}



\begin{quote}
Before the new one: say the Mandarin for to break the traffic rules, then the Mandarin for a main road between towns.
\end{quote}

\subsection*{You'll want to know: \zh{高速}}
\textbf{\zh{高速}} — \emph{gāosù} — \textquotedblleft{}expressway\textquotedblright{}.

\begin{quote}
\textbf{\zh{去高速怎么走}} \emph{qù gāosù zěnme zǒu} — How do I get onto the expressway?
\end{quote}

The new character, in pieces: \zh{速} — ten strokes in two pieces: the right piece first, seven strokes, a bar, a small \zh{口}, a long vertical down through them, a stroke falling to the left and a dot; then the walking piece of \zh{这} under it, three strokes. The next lesson puts it in your hand.

\begin{cousinweb}
\textbf{\zh{高速}} (\emph{gāosù}) is built from \zh{高} 'high' and \zh{速} 'speed', both inherited from Middle Chinese, \emph{kaw} and \emph{suwk} (Guangyun, in Baxter's transcription), which glosses the second 'swift'; it is a short form of a longer name, the high-speed road.
\end{cousinweb}

\begin{grammarlens}[title={What you've built}]
One more word for getting around.
\end{grammarlens}

\subsection*{Your turn}
\begin{itemize}
\raggedright
  \item \emph{Say it:} \emph{gāosù}
  \item \emph{Say it:} \emph{gāosù}, once more
  \item \emph{Say it:} \emph{qù gāosù zěnme zǒu} — How do I get onto the expressway?
  \item \emph{Recall:} say the Mandarin for to break the traffic rules, then the Mandarin for a main road between towns, then say \emph{gāosù} again
\end{itemize}

\begin{tcolorbox}[breakable,colback=teal!4,colframe=teal!35!black,title={Before you move on}]
What does \zh{高速} mean? (\textquotedblleft{}Expressway\textquotedblright{}.) Say it once more.
\end{tcolorbox}
\section[\texorpdfstring{\zh{速} (sù)}{速 (sù)}]{\zh{速} — a bundle, then the walking piece}

\label{lesson:ZH-W421-su-speed}



\begin{quote}
The word first, then the character inside it.

\begin{itemize}
\raggedright
  \item \emph{Recall:} say \emph{gāosù} — \textbf{R1}, one lesson back, before it decays
\end{itemize}

\begin{quote}
\zh{这} \zh{高}
\end{quote}

\emph{Read it:} each one aloud, with nothing said first — \emph{zhè}, \emph{gāo}
\end{quote}

\subsection*{Script — a bundle, then the walking piece}
\hlblockfigure{\detokenize{figures/ZH-W421-su-speed-filmstrip.pdf}}{How \zh{速} is written, stroke by stroke}

\begin{quote}
\zh{速}
\end{quote}

\textbf{Two pieces, and the walking piece comes last.} Count the pieces, not the strokes.

\textbf{The upper piece, seven strokes.} A bar across the top. A small \zh{口} under it, three strokes. A long vertical down through the bar and \zh{口}. Then a stroke falling to the left from the vertical and a dot to the right, as at the foot of the tree piece.

\textbf{The walking piece, three strokes}, as in \zh{这}: a dot, a stroke that zigzags down, then a long sweep along the foot.

7 + 3 = ten.

\subsection*{Your turn}
\begin{itemize}
\raggedright
  \item \emph{Write it:} \zh{速}, all ten, with the model in front of you
  \item \emph{Cover:} write it again from memory, and count the pieces, not the strokes
\end{itemize}

\begin{tcolorbox}[breakable,colback=teal!4,colframe=teal!35!black,title={Before you move on}]
How many strokes in \zh{速}? (\textbf{Ten}.) What is written near the top? (\textbf{A small \zh{口}}.) Which piece comes last? (\textbf{The walking piece}.)

Source: \href{https://raw.githubusercontent.com/chanind/hanzi-writer-data/68d10a4b21150cae5e1ebbd223eed289cf32d90c/data/\%E9\%80\%9F.json}{Hanzi Writer Data \zh{速}, pinned stroke paths and medians}.
\end{tcolorbox}
\section[\texorpdfstring{\zh{车速} (chēsù)}{车速 (chēsù)}]{\zh{车速} (chēsù) — speed of a vehicle}

\label{lesson:ZH-C421-chesu}



\begin{quote}
Before the new one: say the Mandarin for a main road between towns, then the Mandarin for the expressway.
\end{quote}

\subsection*{You'll want to know: \zh{车速}}
\textbf{\zh{车速}} — \emph{chēsù} — \textquotedblleft{}speed of a vehicle\textquotedblright{}.

\begin{quote}
\textbf{\zh{高速的车速很快}} \emph{gāosù de chēsù hěn kuài} — Cars go fast on the expressway.
\end{quote}

\begin{cousinweb}
\textbf{\zh{车速}} (\emph{chēsù}) is built from \zh{车} 'vehicle' and \zh{速} 'speed', both inherited from Middle Chinese, \emph{tsyhae} and \emph{suwk} (Guangyun, in Baxter's transcription).
\end{cousinweb}

\begin{grammarlens}[title={What you've built}]
One more word for getting around.
\end{grammarlens}

\subsection*{Your turn}
\begin{itemize}
\raggedright
  \item \emph{Say it:} \emph{chēsù}
  \item \emph{Say it:} \emph{chēsù}, once more
  \item \emph{Say it:} \emph{gāosù de chēsù hěn kuài} — Cars go fast on the expressway.
  \item \emph{Recall:} say the Mandarin for a main road between towns, then the Mandarin for the expressway, then say \emph{chēsù} again
\end{itemize}

\begin{tcolorbox}[breakable,colback=teal!4,colframe=teal!35!black,title={Before you move on}]
What does \zh{车速} mean? (\textquotedblleft{}Speed of a vehicle\textquotedblright{}.) Say it once more.
\end{tcolorbox}
\section[shōufèizhàn]{shōufèizhàn — toll gate}

\label{lesson:ZH-C421-shoufeizhan}



\begin{quote}
Before the new one: say the Mandarin for the expressway, then the Mandarin for the speed of a vehicle.
\end{quote}

\subsection*{You'll want to know: shōufèizhàn}
\emph{shōufèizhàn} — \textquotedblleft{}toll gate\textquotedblright{}.

\begin{quote}
\emph{wǎng qián kāi jiù shì shōufèizhàn} — Drive straight on and you come to the toll gate.
\end{quote}

Its characters come later; for now it is a word for the ear.

\begin{cousinweb}
\emph{shōufèizhàn} is built from \emph{shōu} 'to collect', \emph{fèi} 'a fee' and \zh{站} 'a station', all inherited from Middle Chinese, \emph{syuw}, \emph{phj+jH} and \emph{treamH} (Guangyun, in Baxter's transcription), which glosses the first 'to gather in'.
\end{cousinweb}

\begin{grammarlens}[title={What you've built}]
One more word for getting around.
\end{grammarlens}

\subsection*{Your turn}
\begin{itemize}
\raggedright
  \item \emph{Say it:} \emph{shōufèizhàn}
  \item \emph{Say it:} \emph{shōufèizhàn}, once more
  \item \emph{Say it:} \emph{wǎng qián kāi jiù shì shōufèizhàn} — Drive straight on and you come to the toll gate.
  \item \emph{Recall:} say the Mandarin for the expressway, then the Mandarin for the speed of a vehicle, then say \emph{shōufèizhàn} again
\end{itemize}

\begin{tcolorbox}[breakable,colback=teal!4,colframe=teal!35!black,title={Before you move on}]
What does shōufèizhàn mean? (\textquotedblleft{}Toll gate\textquotedblright{}.) Say it once more.
\end{tcolorbox}
\section[chāosù]{chāosù — to break the speed limit}

\label{lesson:ZH-C421-chaosu}



\begin{quote}
Before the new one: say the Mandarin for the speed of a vehicle, then the Mandarin for a toll gate.
\end{quote}

\subsection*{You'll want to know: chāosù}
\emph{chāosù} — \textquotedblleft{}to break the speed limit\textquotedblright{}.

\begin{quote}
\emph{nǐ kāi de tài kuài le chāosù le} — You're driving too fast: you're over the limit.
\end{quote}

Its characters come later; for now it is a word for the ear.

\begin{cousinweb}
\emph{chāosù} is built from \emph{chāo} 'to go beyond' and \zh{速} 'speed', both inherited from Middle Chinese, \emph{trhjew} and \emph{suwk} (Guangyun, in Baxter's transcription), which glosses the first 'to leap'.
\end{cousinweb}

\begin{grammarlens}[title={What you've built}]
One more word for getting around.
\end{grammarlens}

\subsection*{Your turn}
\begin{itemize}
\raggedright
  \item \emph{Say it:} \emph{chāosù}
  \item \emph{Say it:} \emph{chāosù}, once more
  \item \emph{Say it:} \emph{nǐ kāi de tài kuài le chāosù le} — You're driving too fast: you're over the limit.
  \item \emph{Recall:} say the Mandarin for the speed of a vehicle, then the Mandarin for a toll gate, then say \emph{chāosù} again
\end{itemize}

\begin{tcolorbox}[breakable,colback=teal!4,colframe=teal!35!black,title={Before you move on}]
What does chāosù mean? (\textquotedblleft{}To break the speed limit\textquotedblright{}.) Say it once more.
\end{tcolorbox}
